% GENERATED FILE. Edit canonical lessons, then run npm run generate:books.
% canonical-source-hash: fnv1a64:b8928ec0620bd4ed
% canonical-lessons: MR-C345-niskarsa, MR-C345-anuman-kadhne, MR-C345-adhar, MR-C345-sarams, MR-C345-thodkyat

\chapter{Reaching Conclusions}
\label{ch:mr-reaching-conclusions}
\hlchaptermodality{345}

\begin{chapteropening}
\textbf{By the end of this chapter:} \emph{I can reach a conclusion and sum things up.}

Complete the last lesson of chapter 345: I can reach a conclusion and sum things up.
\end{chapteropening}

\section[\texorpdfstring{\mr{निष्कर्ष} (niṣkarṣa)}{निष्कर्ष (niṣkarṣa)}]{\mr{निष्कर्ष} (niṣkarṣa) — a conclusion}

\label{lesson:MR-C345-niskarsa}



\begin{quote}
Before the new one: say the Marathi for an inconsistency, then the Marathi for a standpoint.
\end{quote}

\subsection*{You'll want to know: \mr{निष्कर्ष}}
\textbf{\mr{निष्कर्ष}} — \emph{niṣkarṣa} — \textquotedblleft{}a conclusion\textquotedblright{}.

\begin{quote}
\mr{शेवटी} \mr{आम्ही} \mr{एका} \mr{निष्कर्षावर} \mr{आलो}.
\end{quote}

\begin{quote}
\emph{śevṭī āmhī ekā niṣkarṣāvar ālo.}
\end{quote}

\begin{quote}
In the end we came to a conclusion.
\end{quote}

Masculine.

\mr{निष्कर्ष} \mr{काढणे}, to draw a conclusion.

\begin{cousinweb}
\textbf{\mr{निष्कर्ष}} (\emph{niṣkarṣa}) is a tatsama, taken from Sanskrit \emph{niṣkarṣa}, \textquotedblleft{}drawing out, extracting\textquotedblright{} (Monier-Williams). A conclusion is what is drawn out of the facts.
\end{cousinweb}

\begin{grammarlens}[title={What you've built}]
One more word for giving reasons and opinions.
\end{grammarlens}

\subsection*{Your turn}
\begin{itemize}
\raggedright
  \item \emph{Say it:} \emph{niṣkarṣa}
  \item \emph{Say it:} \emph{niṣkarṣa}, once more
  \item \emph{Say it:} \emph{śevṭī āmhī ekā niṣkarṣāvar ālo.} — In the end we came to a conclusion.
  \item \emph{Recall:} say the Marathi for an inconsistency, then the Marathi for a standpoint, then say \emph{niṣkarṣa} again
\end{itemize}

\begin{tcolorbox}[breakable,colback=teal!4,colframe=teal!35!black,title={Before you move on}]
What does \mr{निष्कर्ष} mean? (\textquotedblleft{}A conclusion\textquotedblright{}.) Say it once more.
\end{tcolorbox}
\section[\texorpdfstring{\mr{अनुमान} \mr{काढणे} (anumān kāḍhṇe)}{अनुमान काढणे (anumān kāḍhṇe)}]{\mr{अनुमान} \mr{काढणे} (anumān kāḍhṇe) — to infer}

\label{lesson:MR-C345-anuman-kadhne}



\begin{quote}
Before the new one: say the Marathi for a standpoint, then the Marathi for a conclusion.
\end{quote}

\subsection*{You'll want to know: \mr{अनुमान} \mr{काढणे}}
\textbf{\mr{अनुमान} \mr{काढणे}} — \emph{anumān kāḍhṇe} — \textquotedblleft{}to infer\textquotedblright{}.

\begin{quote}
\mr{काळं} \mr{आभाळ} \mr{पाहून} \mr{त्याने} \mr{पावसाचं} \mr{अनुमान} \mr{काढलं}.
\end{quote}

\begin{quote}
\emph{kāḷaṁ ābhāḷ pāhūn tyāne pāvsācaṁ anumān kāḍhlaṁ.}
\end{quote}

\begin{quote}
Seeing the dark sky, he inferred that it would rain.
\end{quote}

Past, with -\mr{ने} (he …): \mr{त्याने} \mr{अनुमान} \mr{काढलं}, \emph{tyāne anumān kāḍhlaṁ}.

\mr{अनुमान}, an inference; \mr{काढणे}, to draw out.

\begin{cousinweb}
\textbf{\mr{अनुमान} \mr{काढणे}} (\emph{anumān kāḍhṇe}) joins the tatsama \mr{अनुमान}, Sanskrit \emph{anumāna}, \textquotedblleft{}inference, thought\textquotedblright{} (CDIAL 328), to \mr{काढणे}, \textquotedblleft{}to take out\textquotedblright{}.
\end{cousinweb}

\begin{grammarlens}[title={What you've built}]
One more word for giving reasons and opinions.
\end{grammarlens}

\subsection*{Your turn}
\begin{itemize}
\raggedright
  \item \emph{Say it:} \emph{anumān kāḍhṇe}
  \item \emph{Say it:} \emph{anumān kāḍhṇe}, once more
  \item \emph{Say it:} \emph{kāḷaṁ ābhāḷ pāhūn tyāne pāvsācaṁ anumān kāḍhlaṁ.} — Seeing the dark sky, he inferred that it would rain.
  \item \emph{Recall:} say the Marathi for a standpoint, then the Marathi for a conclusion, then say \emph{anumān kāḍhṇe} again
\end{itemize}

\begin{tcolorbox}[breakable,colback=teal!4,colframe=teal!35!black,title={Before you move on}]
What does \mr{अनुमान} \mr{काढणे} mean? (\textquotedblleft{}To infer\textquotedblright{}.) Say it once more.
\end{tcolorbox}
\section[\texorpdfstring{\mr{आधार} (ādhār)}{आधार (ādhār)}]{\mr{आधार} (ādhār) — a basis; support}

\label{lesson:MR-C345-adhar}



\begin{quote}
Before the new one: say the Marathi for a conclusion, then the Marathi for to infer.
\end{quote}

\subsection*{You'll want to know: \mr{आधार}}
\textbf{\mr{आधार}} — \emph{ādhār} — \textquotedblleft{}a basis; support\textquotedblright{}.

\begin{quote}
\mr{तुमच्या} \mr{मताला} \mr{काय} \mr{आधार} \mr{आहे}?
\end{quote}

\begin{quote}
\emph{tumacyā matālā kāy ādhār āhe?}
\end{quote}

\begin{quote}
What basis does your view have?
\end{quote}

Masculine.

X \mr{ला} \mr{आधार} \mr{देणे}, to support X; \mr{च्या} \mr{आधारावर}, on the basis of.

\begin{cousinweb}
\textbf{\mr{आधार}} (\emph{ādhār}) is a tatsama, taken from Sanskrit \emph{ādhāra}, \textquotedblleft{}support, base\textquotedblright{} (CDIAL 1165); the same entry lists an inherited Marathi \emph{āhārā}, a ring of grass that holds a pot steady.
\end{cousinweb}

\begin{grammarlens}[title={What you've built}]
One more word for giving reasons and opinions.
\end{grammarlens}

\subsection*{Your turn}
\begin{itemize}
\raggedright
  \item \emph{Say it:} \emph{ādhār}
  \item \emph{Say it:} \emph{ādhār}, once more
  \item \emph{Say it:} \emph{tumacyā matālā kāy ādhār āhe?} — What basis does your view have?
  \item \emph{Recall:} say the Marathi for a conclusion, then the Marathi for to infer, then say \emph{ādhār} again
\end{itemize}

\begin{tcolorbox}[breakable,colback=teal!4,colframe=teal!35!black,title={Before you move on}]
What does \mr{आधार} mean? (\textquotedblleft{}A basis; support\textquotedblright{}.) Say it once more.
\end{tcolorbox}
\section[\texorpdfstring{\mr{सारांश} (sārāṁś)}{सारांश (sārāṁś)}]{\mr{सारांश} (sārāṁś) — a summary}

\label{lesson:MR-C345-sarams}



\begin{quote}
Before the new one: say the Marathi for to infer, then the Marathi for a basis.
\end{quote}

\subsection*{You'll want to know: \mr{सारांश}}
\textbf{\mr{सारांश}} — \emph{sārāṁś} — \textquotedblleft{}a summary\textquotedblright{}.

\begin{quote}
\mr{मीराने} \mr{धड्याचा} \mr{सारांश} \mr{सांगितला}.
\end{quote}

\begin{quote}
\emph{mīrāne dhaḍyācā sārāṁś sāṅgitlā.}
\end{quote}

\begin{quote}
Meera gave a summary of the lesson.
\end{quote}

Masculine.

X \mr{चा} \mr{सारांश}, a summary of X.

\begin{cousinweb}
\textbf{\mr{सारांश}} (\emph{sārāṁś}) is a learned compound built from Sanskrit \emph{sāra}, \textquotedblleft{}essence, best part\textquotedblright{}, and \emph{aṁśa}, \textquotedblleft{}part\textquotedblright{} (Monier-Williams): the essential part.
\end{cousinweb}

\begin{grammarlens}[title={What you've built}]
One more word for giving reasons and opinions.
\end{grammarlens}

\subsection*{Your turn}
\begin{itemize}
\raggedright
  \item \emph{Say it:} \emph{sārāṁś}
  \item \emph{Say it:} \emph{sārāṁś}, once more
  \item \emph{Say it:} \emph{mīrāne dhaḍyācā sārāṁś sāṅgitlā.} — Meera gave a summary of the lesson.
  \item \emph{Recall:} say the Marathi for to infer, then the Marathi for a basis, then say \emph{sārāṁś} again
\end{itemize}

\begin{tcolorbox}[breakable,colback=teal!4,colframe=teal!35!black,title={Before you move on}]
What does \mr{सारांश} mean? (\textquotedblleft{}A summary\textquotedblright{}.) Say it once more.
\end{tcolorbox}
\section[\texorpdfstring{\mr{थोडक्यात} (thoḍkyāt)}{थोडक्यात (thoḍkyāt)}]{\mr{थोडक्यात} (thoḍkyāt) — in short; briefly}

\label{lesson:MR-C345-thodkyat}



\begin{quote}
Before the new one: say the Marathi for a basis, then the Marathi for a summary.
\end{quote}

\subsection*{You'll want to know: \mr{थोडक्यात}}
\textbf{\mr{थोडक्यात}} — \emph{thoḍkyāt} — \textquotedblleft{}in short; briefly\textquotedblright{}.

\begin{quote}
\mr{थोडक्यात} \mr{सांगायचं} \mr{तर}, \mr{आम्हाला} \mr{उशीर} \mr{झाला}.
\end{quote}

\begin{quote}
\emph{thoḍkyāt sāṅgāycaṁ tar, āmhālā uśīr jhālā.}
\end{quote}

\begin{quote}
To put it briefly, we were late.
\end{quote}

\mr{थोडक्यात} \mr{सांगायचं} \mr{तर}, to put it briefly; also narrowly: \mr{थोडक्यात} \mr{वाचला}, had a narrow escape.

\begin{cousinweb}
\textbf{\mr{थोडक्यात}} (\emph{thoḍkyāt}) is built from \mr{थोडका}, a form of \mr{थोडा}, \textquotedblleft{}a little\textquotedblright{}, with the ending -\mr{त}, \textquotedblleft{}in\textquotedblright{}: in a little space. \mr{थोडा} is inherited from Sanskrit \emph{stoka}, \textquotedblleft{}little\textquotedblright{} (CDIAL 13720).
\end{cousinweb}

\begin{grammarlens}[title={What you've built}]
One more word for giving reasons and opinions.
\end{grammarlens}

\subsection*{Your turn}
\begin{itemize}
\raggedright
  \item \emph{Say it:} \emph{thoḍkyāt}
  \item \emph{Say it:} \emph{thoḍkyāt}, once more
  \item \emph{Say it:} \emph{thoḍkyāt sāṅgāycaṁ tar, āmhālā uśīr jhālā.} — To put it briefly, we were late.
  \item \emph{Recall:} say the Marathi for a basis, then the Marathi for a summary, then say \emph{thoḍkyāt} again
\end{itemize}

\begin{tcolorbox}[breakable,colback=teal!4,colframe=teal!35!black,title={Before you move on}]
What does \mr{थोडक्यात} mean? (\textquotedblleft{}In short; briefly\textquotedblright{}.) Say it once more.
\end{tcolorbox}
